% 黄道十二宫（综述）
% license CCBYSA3
% type Wiki

本文根据 CC-BY-SA 协议转载翻译自维基百科\href{https://en.wikipedia.org/wiki/Astrological_sign}{相关文章}。

\begin{figure}[ht]
\centering
\includegraphics[width=6cm]{./figures/e66f90583cecd88b.png}
\caption{出自《贝里公爵的豪华时祷书》的解剖学黄道人} \label{fig_Astsi_1}
\end{figure}
在西方占星术中，占星星座即黄道，被划分为十二个各占30度的区段，这些区段由太阳在地球天空中所观测到的360度轨道路径所穿越。星座的顺序从春季的第一天开始，即被称为“白羊宫起点”的位置，也就是春分点。占星星座依次为白羊座、金牛座、双子座、巨蟹座、狮子座、处女座、天秤座、天蝎座、人马座、摩羯座、宝瓶座和双鱼座。西方黄道起源于巴比伦占星术，后来受到希腊化文化的影响。每一个星座均以太阳每年穿越天空时经过的一个星座命名。这一观测在简化且流行的太阳星座占星术中被尤为强调。几个世纪以来，由于地球自转轴进动$^{[1]}$，西方占星术中的黄道分区逐渐与其所命名的星座发生偏离，而印度占星术的测量方法则对这种偏移进行了修正。$^{[2]}$ 占星术（即一种基于天象的预兆体系）在中国和西藏文化中也有所发展，但这些占星体系并不以黄道为基础，而是涉及整个天空。

占星术是一种伪科学。$^{[3]}$ 对其理论基础$^{[4]}$以及对相关主张的实验验证$^{[5]}$的科学研究表明，它既不具有科学有效性，也缺乏解释力。对于人格特征与出生月份之间表面相关性的现象，存在更为合理的解释，例如人类季节性出生所产生的影响。

根据占星学的观点，天体现象与人类活动之间遵循“上如其下”的原则，因此星座被认为代表着具有特征性的表达模式。$^{[6]}$ 直到19世纪之前，科学天文学与西方占星术一样，均使用相同的黄道分区。

目前，不同的占星体系采用了多种测量和划分天空的方法，尽管黄道星座的名称与符号传统在很大程度上保持一致。西方占星术以春分点和至点（即热带年中昼夜等长、最长与最短的时刻）为基准进行测量，而印度占星术则沿着赤道平面进行测量，对应的是恒星年。
\subsection{西方黄道星座}  
\subsubsection{历史}
\begin{figure}[ht]
\centering
\includegraphics[width=6cm]{./figures/c64ed924a06babf0.png}
\caption{十二个黄道星座。每一个圆点标示一个星座的起点，相邻星座之间相隔30°。天赤道与黄道的交点定义了分点：白羊宫起点（）与天秤宫起点（）。包含天极与黄极的大圆（P 与 P′）在黄道上分别与 0° 巨蟹宫（）和 0° 摩羯宫（）相交。在该示意图中，太阳被示意性地放置在宝瓶宫起点（）。} \label{fig_Astsi_2}
\end{figure}
西方占星术是希腊化占星术的直接延续，后者在公元2世纪由托勒密记录于《四书》（Tetrabiblos）之中。希腊化占星术本身又部分建立在巴比伦传统的概念之上。具体而言，将黄道划分为十二个相等区段这一思想结构源自巴比伦。$^{[7]}$ 这种黄道划分起源于古代汇编文献《MUL.APIN》中所记载的巴比伦“理想历法”，以及其与巴比伦阴历的结合，后者在《MUL.APIN》中被表述为“月亮之路”。$^{[8]}$ 从某种意义上说，黄道正是对一种理想化阴历体系的抽象化表达。

至公元前4世纪，巴比伦天文学及其天象征兆体系对古希腊文化产生了影响；而到公元前2世纪末，埃及天文学同样发挥了重要作用。这一发展不同于美索不达米亚传统，其结果是占星学的重心转向对个体出生星图的关注，并催生了本命占星术（Horoscopic astrology），引入了上升点（即出生时刻黄道升起的度数）以及十二宫的概念。将占星星座与恩培多克勒的四种古典元素相联系，也是刻画十二星座特征的另一项重要发展。

至公元2世纪时，希腊化占星传统的整体体系由托勒密在《四书》中加以系统阐述。该著作不仅是后世西方天文学传统的奠基性文本，也对印度和伊斯兰世界产生了深远影响，并在近十七个世纪中一直作为权威参考，因为后续传统几乎未对其核心教义作出实质性修改。
\subsubsection{西方占星对应表}  
下表展示了十二个占星星座的大致日期范围，以及每个星座在古典体系$^{[9]}$和现代体系$^{[10]}$中的守护星。按照定义，白羊宫始于白羊宫起点，即太阳位于三月春分点时的位置。春分的具体日期每年略有变化，但始终介于3月19日至3月21日之间。因此，白羊宫的起始日期，以及其他所有星座的起始日期，都会随年份略有变动。下述西方占星表列举了天球经度的十二个分区，并使用拉丁名称表示。经度区间在定义上，对第一个端点（a）为闭区间，对第二个端点（b）为开区间——例如，经度30°是金牛宫的起点，而不属于白羊宫。在天文学文献中，这些星座有时会以0至11的编号代替其符号来表示。

